\subsection{トピックと参考文献の対応}
この表は、SWEBOKが示すトピックと参考資料の対応を、学習用に簡略化したものである。離散数学の基礎にはRosenを、数値計算にはCheney and Kincaidを読むとよい。

\par\smallskip\noindent\refstepcounter{table}\label{tab:ch17-22}表 \thetable: トピックと参考文献の対応におけるトピック・主な参考資料\par\smallskip
表 \thetable は、トピックと参考文献の対応におけるトピック・主な参考資料を比較・確認しやすい形で整理したものである。\par\smallskip
\begin{longtable}{>{\raggedright\arraybackslash}p{0.26\linewidth} >{\raggedright\arraybackslash}p{0.68\linewidth}}
\toprule
トピック & 主な参考資料 \\
\midrule
\endfirsthead
\toprule
トピック & 主な参考資料 \\
\midrule
\endhead
1. 基本論理 & Rosen, Discrete Mathematics and Its Applications \\
2. 証明技法 & Rosen \\
3. 集合・関係・関数 & Rosen \\
4. グラフと木 & Rosen \\
5. 有限状態機械 & Rosen \\
6. 文法 & Rosen \\
7. 数論 & Rosen \\
8. 数え上げの基礎 & Rosen \\
9. 離散確率 & Rosen \\
10. 数値精度・正確度・誤差 & Cheney and Kincaid, Numerical Mathematics and Computing \\
11. 代数構造 & Rosen \\
12. 工学微積分 & Cheney and Kincaid \\
\bottomrule
\end{longtable}
